Throughout $e_{a}e_{b} + e_{b}e_{a} = -2\delta_{ab}$, vectors act on $S$ by left multiplication, and $\bar{a}$ is the Clifford conjugate (reversal composed with the grade involution), for which $\bar{B}B = 1$ on every basis blade $B$ of $\Cl(0,6)$. So $(a, b) = \langle\bar{a}b\rangle_{0}$ is a positive definite scalar product on $\Cl(0,6)$ and on $S$, and left multiplication by a vector $v$ is skew for it, $(va, b) = -(a, vb)$, since $\bar{v} = -v$; left multiplication by a bivector is skew as well. The check builds $S$ as $\Cl(0,6)P$ with $P = \tfrac18(1 + e_{1}e_{2}e_{3})(1 + e_{1}e_{4}e_{5})(1 + e_{2}e_{4}e_{6})$, three commuting blades of square $+1$, and verifies $P^{2} = P$, $\dim S = 8$, $J^{2} = -1$, and that $J$ anticommutes with vectors and commutes with bivectors.

\emph{(i).} On $\mathcal{P}_{k}$ put the Fischer product $[P, R] = \sum_{\alpha}\alpha!\,(p_{\alpha}, r_{\alpha})$ for $P = \sum p_{\alpha}x^{\alpha}$, $R = \sum r_{\alpha}x^{\alpha}$. Then $[x_{a}P, R] = [P, \partial_{a}R]$ for $P \in \mathcal{P}_{k-1}$, $R \in \mathcal{P}_{k}$, and with the skewness of $e_{a}$,
\[
[xP, R] = \sum_{a}[e_{a}x_{a}P, R] = -\sum_{a}[x_{a}P, e_{a}R] = -[P, \partial R] .
\]
So multiplication by $x$ is minus the adjoint of $\partial$. It is injective, since $x(xP) = Q(x)P = -|x|^{2}P$. Hence the orthogonal complement of $\partial\mathcal{P}_{k}$ in $\mathcal{P}_{k-1}$, which is the kernel of $x$ on $\mathcal{P}_{k-1}$, is zero: $\partial$ is onto, $\dim M_{k} = \dim\mathcal{P}_{k} - \dim\mathcal{P}_{k-1} = 8[\binom{k+5}{5} - \binom{k+4}{5}] = 8\binom{k+4}{4}$, and $\mathcal{P}_{k} = M_{k} \oplus x\mathcal{P}_{k-1}$, the second summand being the orthogonal complement of the kernel of $\partial$. For $P \in M_{k}$,
\[
\partial(xP) = \sum_{a}e_{a}e_{a}P + \sum_{a,b}x_{b}\,e_{a}e_{b}\,\partial_{a}P = -6P + \sum_{a,b}x_{b}(-e_{b}e_{a} - 2\delta_{ab})\partial_{a}P = -6P - x\,\partial P - 2EP = -(6 + 2k)P .
\]
Iterating the decomposition, $\mathcal{P}_{k} = \bigoplus_{j}x^{j}M_{k-j}$, and on the sphere $x^{2} = -1$, so $x^{2j}M_{k-2j}$ and $x^{2j+1}M_{k-2j-1}$ restrict to $\pm M_{k-2j}$ and $\pm xM_{k-2j-1}$. The check computes $M_{k}$ as an exact kernel for $k \le 3$: dimensions $8, 40, 120, 280$.

\emph{(ii).} $x\partial = \sum_{a,b}x_{a}e_{a}e_{b}\partial_{b} = -E + x \wedge \partial$, the terms $a = b$ giving $-E$ and the others the bivector part. On $M_{k}$, $0 = x\partial P = -kP + (x \wedge \partial)P$. On $xP$, $P \in M_{k}$, by (i) $x\partial(xP) = -(6 + 2k)xP$ and $E(xP) = (k+1)xP$, so $(x \wedge \partial)(xP) = -(6 + 2k)xP + (k + 1)xP = -(k + 5)xP$.

\emph{(iii).} For $X$, $Y$ tangent at $n$, $Xn = -nX$ and $n^{2} = -1$, so $(Xn)(Yn) = -XYn^{2} = XY$ and the Clifford relation follows; $(Xn)^{2} = Q(X)$. The Levi-Civita derivative of the sphere is the tangential part of the flat one: $\partial_{X}Y = D_{X}Y - \langle X, Y\rangle n$ with $\langle X, Y\rangle = -\tfrac12(XY + YX)$, and $\partial_{X}n = X$. Then
\[
\nabla_{X}(Yn\,\psi) - (D_{X}Y)n\,\psi - Yn\,\nabla_{X}\psi = \bigl(-\langle X, Y\rangle n^{2} + YX + \tfrac12(XnYn - YnXn)\bigr)\psi = \bigl(\langle X, Y\rangle + YX + \tfrac12(XY - YX)\bigr)\psi = 0 .
\]
The term $\tfrac12Xn$ is a bivector, skew for $(\cdot, \cdot)$, so $\nabla$ preserves the scalar product, and it commutes with $J$. It is the connection that the flat connection of the six-space induces on the hypersurface, the lift of the Levi-Civita connection for the spin structure of the sphere (which is unique, the sphere being simply connected); this identification is the standard hypersurface formula and is not needed for the computation. In a tangent frame, $\sum_{a}e_{a}n\,\partial_{a} = -n\sum_{a}e_{a}\partial_{a} = -n(\partial - n\partial_{n}) = -x\partial - \partial_{n}$ on the sphere, with $\partial_{n}$ the radial derivative, equal to $E$ there; so $\sum_{a}e_{a}n\,\partial_{a} = E - x \wedge \partial - E = -x \wedge \partial$, an operator that differentiates along the sphere only. And $\tfrac12\sum_{a}(e_{a}n)(e_{a}n) = \tfrac12\sum_{a}Q(e_{a}) = -\tfrac52$. Hence $D_{\Sph} = -x \wedge \partial - \tfrac52$, and by (ii) its values on $M_{k}$ and $xM_{k}$ are $-(k + \tfrac52)$ and $k + 5 - \tfrac52 = k + \tfrac52$. $D_{\Sph}$ is symmetric for the integrated scalar product, since $\nabla$ is metric and the Clifford map is skew, so restrictions with different eigenvalues are orthogonal, and those with the same eigenvalue have the dimension of $M_{k}$ because restriction of a homogeneous polynomial to the sphere is injective. Completeness: by (i) the eigenfunctions span the restrictions of all $S$-valued polynomials, which are dense in the continuous $S$-valued functions on the sphere (Weierstrass), hence in the square-integrable ones. So the spectrum is $\pm(k + \tfrac52)$ with real multiplicity $8\binom{k+4}{4}$, that is $4\binom{k+4}{4}$ over $\mathrm{span}\{1, J\}$, the classical spectrum of the round five-sphere.

\emph{(iv).} The antipodal map $a(x) = -x$ has differential $X \mapsto -X$ and carries $n$ to $-n$, so it preserves $Xn$; $J$ commutes with bivectors, hence with $Xn$, with $\tfrac12Xn$ and with $x \wedge \partial$, and $x \wedge \partial$ is unchanged by $x \mapsto -x$. So $A = J \circ a^{*}$ commutes with $\nabla$ and $D_{\Sph}$, and $A^{2} = J^{2} = -1$. A homogeneous polynomial of degree $k$ has parity $(-1)^{k}$, so $A = (-1)^{k}J$ on $M_{k}$ and $(-1)^{k+1}J$ on $xM_{k}$, where $J(xP) = -x(JP)$ because $J$ anticommutes with vectors; $JP \in M_{k}$ since $\partial(JP) = -J\partial P$. The antipodal map is the rotation $-1$ of the six-space, the product of the half turns in the planes $e_{1}e_{2}$, $e_{3}e_{4}$, $e_{5}e_{6}$; the rotors covering it are $\pm J = \pm(e_{1}e_{2})(e_{3}e_{4})(e_{5}e_{6})$, each bivector $e_{a}e_{b}$ being the rotor of a half turn, and both have square $-1$. A spin structure on the quotient is a lift of the deck group to the rotor group, so there is none. (The operator $\psi(x) \mapsto \psi(-x)$ alone is an involution commuting with the Clifford map, but it is not the action of a rotor; it is $A$ composed with $-J$, an element of order four of $\mathrm{span}\{1, J\}$, which is the pairing made below with $\omega$.) By \claimref{synthesis.order-four-elements-with-square-the-grading-and-the-majorana-element}(2), for the ten Majorana generators of Hypothesis (S) (\claimref{standard-model.hyp-s-slot-statistics}), $\omega^{2} = -1$ and $\omega = -I$ on the sixteen $V_{+}$. On $S \otimes V_{+}$ the operators $J \otimes 1$ and $1 \otimes I$ commute and square to $-1$, so $J \otimes I$ is an involution; $W_{+}$ is its $-1$ eigenspace, where $J \otimes 1 = 1 \otimes I$, and there $J \otimes \omega = -(J \otimes 1)(1 \otimes I) = -(J \otimes 1)^{2} = +1$; on the $+1$ eigenspace $W_{-}$, $J \otimes \omega = -1$. Both halves are preserved by $D_{\Sph} \otimes 1$ and by $A \otimes \omega$, and $W_{+} \cong S \otimes_{I} V_{+}$ is, per component of $V_{+}$, a copy of $S$. On $W_{+}$-valued fields $(A \otimes \omega)\Psi(x) = (J \otimes \omega)\Psi(-x) = \Psi(-x)$. The fields invariant under the deck transformation are the even ones; their eigenvalues are those of (iii) on the $M_{k}$ with $k$ even and the $xM_{k}$ with $k$ odd. The check verifies $(J\omega)^{2} = 1$ in $\Cl(10,6)$, where $J$ and $\omega$ are the volume elements of the two summands, and the parity statements on bases of $M_{k}$, $k \le 3$.

\emph{(v).} For even fields the positive eigenvalues are $k + \tfrac52$, $k$ odd, and the negative ones $-(k + \tfrac52)$, $k$ even, with multiplicity $m_{k} = 4\binom{k+4}{4} = \tfrac16(k+1)(k+2)(k+3)(k+4)$ over the unit, which gives the stated $\eta(s)$ for $\operatorname{Re}s > 5$; with $u = k + \tfrac52$ the four factors are $u \mp \tfrac32$, $u \mp \tfrac12$. The binomial series gives $\sum_{k}(-1)^{k}\binom{k+4}{4}t^{k} = (1 + t)^{-5}$ for $|t| < 1$, with limit $2^{-5}$ at $t \to 1$: Abel sum $4 \cdot 2^{-5} = \tfrac18$. For the continuation write $m_{k} = \sum_{j}c_{j}(k + a)^{j}$ and $F_{a}(s) = \sum_{k \ge 0}(-1)^{k}(k + a)^{-s}$, $a > 0$. For $\operatorname{Re}s > 0$,
\[
\Gamma(s)F_{a}(s) = \int_{0}^{\infty}t^{s-1}f(t)\,dt, \qquad f(t) = \frac{e^{-at}}{1 + e^{-t}},
\]
and $f$ is smooth on $[0, \infty)$ and decays exponentially. For such $f$ the function $\Gamma(s)^{-1}\int_{0}^{\infty}t^{s-1}f(t)\,dt$ continues to an entire function with value $(-1)^{j}f^{(j)}(0)$ at $s = -j$ (integrate by parts $j + 1$ times; the classical Mellin lemma). Now $f(t) = g(-t)$ with $g(z) = e^{az}/(1 + e^{z}) = \tfrac12\sum_{j}E_{j}(a)z^{j}/j!$, the generating function of the Euler polynomials, so $F_{a}(-j) = g^{(j)}(0) = \tfrac12E_{j}(a)$, and
\[
-\eta(0) = \sum_{j}c_{j}F_{a}(-j) = \tfrac12\sum_{j}c_{j}E_{j}(a) .
\]
The same number is the Abel sum: $\sum_{k}(-1)^{k}(k + a)^{j}e^{(k + a)z} = g^{(j)}(z)$ for $\operatorname{Re}z < 0$, with limit $g^{(j)}(0)$ as $z \to 0$, and $\sum_{k}(-1)^{k}m_{k}e^{(k+a)z} = 4e^{az}(1 + e^{z})^{-5}$ tends to $\tfrac18$. With $a = \tfrac52$, $c_{4} = \tfrac16$, $c_{2} = -\tfrac5{12}$, $c_{0} = \tfrac3{32}$, $E_{2}(a) = a^{2} - a = \tfrac{15}{4}$, $E_{4}(a) = a^{4} - 2a^{3} + a = \tfrac{165}{16}$: $\tfrac12(\tfrac16 \cdot \tfrac{165}{16} - \tfrac5{12} \cdot \tfrac{15}{4} + \tfrac3{32}) = \tfrac18$. The check verifies the generating function to order $15$, the Euler polynomials against $2e^{az}/(e^{z} + 1)$, and the value $\tfrac18$ for the shifts $\tfrac52$, $1$, $\tfrac12$ and $\tfrac73$. No eigenvalue vanishes, so $\xi = \tfrac12\eta(0) = -\tfrac1{16}$ for even fields; for odd fields every sign is reversed. Components on which the order-four element acts by the same unit have the same parity of fields and add; a component on which it acts by the opposite unit lies in $W_{+}$ for the opposite identification and has fields of the opposite parity, contributing the opposite sign.
